\section*{Chapter \ref{ch:dv:asians-lookbacks-cliquets-and-forward-starts} --- Asians, Lookbacks, Cliquets and Forward-Starts}

\begin{solution}{exo:dv:asians-lookbacks-cliquets-and-forward-starts:1}
$n=4$: $\sqrt{5\times9/(6\times16)}=\sqrt{0.469}=0.685$. As $n\to\infty$ the factor tends to $\sqrt{2n^2/6n^2}=1/\sqrt3=0.577$.
\end{solution}

\begin{solution}{exo:dv:asians-lookbacks-cliquets-and-forward-starts:2}
Path by path, the arithmetic mean of positive numbers is at least their geometric mean, and the call payoff is increasing
in the average; so the arithmetic payoff dominates the geometric one on every path.
\end{solution}

\begin{solution}{exo:dv:asians-lookbacks-cliquets-and-forward-starts:3}
The six-month at-the-money call today is $S\,C_{\mathrm{BS}}(1,1,0.5)$ by homogeneity; the forward-start call is
$Se^{-q\times0.5}C_{\mathrm{BS}}(1,1,0.5)$. The ratio is $e^{-0.01}=0.990$: the dividends paid before the strike is set lower
the spot at which it is set.
\end{solution}

\begin{solution}{exo:dv:asians-lookbacks-cliquets-and-forward-starts:4}
If $r\le f$: $f+0-0=f$. If $f<r<c$: $f+(r-f)-0=r$. If $r\ge c$: $f+(r-f)-(r-c)=c$. So the clipped return is a constant plus a
call spread between $f$ and $c$, $(r-f)^+-(r-c)^+$. When $c-f$ is small the spread is $(c-f)$ times a digital at about
$(f+c)/2$: each period pays like a forward-start digital, whose price is set by the forward skew.
\end{solution}

\begin{solution}{exo:dv:asians-lookbacks-cliquets-and-forward-starts:5}
\omterm{def:dv:local-volatility:model}{Local volatility} fits today's skew with a volatility that is a fixed function of spot and time; as the forward date
recedes, the spot has moved and the local skew it meets is the one calibrated to decay with maturity, so the \omterm{def:dv:local-volatility:forward}{forward
smile} flattens. In Heston the skew comes from the correlation of the spot with a variance that is still random at the
forward date: it is regenerated there, and the randomness of the variance adds curvature.
\end{solution}

\begin{solution}{exo:dv:asians-lookbacks-cliquets-and-forward-starts:6}
The standard deviation of the plain estimator per path is about $0.025\sqrt{100\,000}=8.0$, so an error of 0.001 needs
$(8.0/0.001)^2=64$ million paths; with the control variate about 47\,000. Computing the geometric payoff adds little to each path,
so the time falls by about the same factor.
\end{solution}

\begin{solution}{exo:dv:asians-lookbacks-cliquets-and-forward-starts:7}
With coupons of 15\%, 25\% and 35\%: \omterm{def:dv:local-volatility:model}{local volatility} 1.23\%, 5.43\% and 12.06\%; Heston 2.49\%, 7.56\% and 14.55\%; Black--Scholes
0.58\%, 3.57\% and 10.00\%. Heston's price moves most in points (12.1 against 10.8 and 9.4); at the low coupon it is twice the
local-volatility price, because it gives more weight to paths with many large falls.
\end{solution}

\begin{solution}{exo:dv:asians-lookbacks-cliquets-and-forward-starts:8}
Vanillas pin the terminal distributions at each expiry, not the joint law of the returns between them. A \omterm{def:dv:asians-lookbacks-cliquets-and-forward-starts:cliquet}{cliquet} pays on
forward returns, so its price depends on the \omterm{def:dv:local-volatility:forward}{forward smile}, which the local-volatility model forecasts from its own
assumptions. Another model fitting the same vanillas prices the chapter's \omterm{def:dv:asians-lookbacks-cliquets-and-forward-starts:cliquet}{cliquet} 45\% higher; the fit is not evidence for
either price.
\end{solution}

\begin{solution}{pb:dv:asians-lookbacks-cliquets-and-forward-starts:1}
\textbf{1.} 5.32 and 8.83.
\textbf{2.} 0.376; $0.61\times20\%=12.3\%$.
\textbf{3.} 0.025 without and 0.0007 with the control variate.
\textbf{4.} 15.69 continuous; daily 15.09 by simulation and 15.08 by the shifted formula.
\textbf{5.} The correction is an expansion in $\sigma\sqrt{\Delta t}$; with a quarter between dates that is 0.1, too large, and
the minimum between dates is far from the monitored one.
\textbf{6.} \omterm{def:dv:local-volatility:model}{Local volatility} 19.19\%, Heston 18.99\%, market 19.19\%.
\textbf{7.} \omterm{def:dv:local-volatility:model}{Local volatility} 21.5\%, Heston 18.5\%; today's one-month 14.9\%.
\textbf{8.} \omterm{def:dv:local-volatility:model}{Local volatility} $23.9-21.3=2.6$ points, Heston $24.0-20.0=4.0$; today's one-month $19.6-12.0=7.6$.
\textbf{9.} The surface's term structure rises: the variance between six and seven months is higher than over the first
month.
\textbf{10.} The one whose forward skew matches forward-start quotes; a \omterm{def:dv:stochastic-volatility:sv}{stochastic volatility model} regenerates the skew,
which \omterm{def:dv:local-volatility:model}{local volatility} does not, though Bergomi found Heston likely to overweight low-volatility, high-skew scenarios
for short forward-starts.
\textbf{11.} 1.53\%, 2.22\% and 1.19\% of notional.
\textbf{12.} At $\pm0.5\%$: 0.78\%, 1.16\% and 0.61\%. At $\pm5\%$: 5.73\%, 6.12\% and 4.74\%.
\textbf{13.} A wide cap makes each period's payoff nearly the return itself, linear, whose value does not depend on the
forward skew; the price then reflects the forward volatility level and the global floor, on which the models differ less.
\textbf{14.} 5.43\%, 7.56\% and 3.57\%.
\textbf{15.} The seller is short forward puts and would prefer the cheaper local-volatility price; that is the warning,
because the model that makes the product cheap is the one whose \omterm{def:dv:local-volatility:forward}{forward smile} has lost its skew.
\textbf{16.} \omterm{def:dv:asians-lookbacks-cliquets-and-forward-starts:fwdstart}{Forward-start options} and \omterm{def:dv:asians-lookbacks-cliquets-and-forward-starts:cliquet}{cliquets} themselves when quoted; options on the volatility index, which reveal the
forward volatility's distribution; \omterm{def:dv:rough-volatility-and-forward-variance-models:fwdvar}{forward variance} from \omterm{def:dv:variance-swaps-and-volatility-derivatives:vs}{variance swaps}.
\textbf{17.} With forward-start risk reversals where they trade; otherwise with vanilla risk reversals of two expiries,
sized by the model, and a reserve for the model's error.
\textbf{18.} At least a large fraction of the gap to the alternative model: 0.69 point of notional on this trade.
\textbf{19.} 0.69 point of notional: 2.22\% under Heston against 1.53\% under \omterm{def:dv:local-volatility:model}{local volatility}, 45\% more, with both fitted to
the same surface.
\textbf{20.} The \omterm{def:dv:local-volatility:forward}{forward smile}: how the smile will look at future dates, which no vanilla pays on.
\end{solution}

\begin{solution}{iq:dv:asians-lookbacks-cliquets-and-forward-starts:1}
The average moves less than the final price: its variance is about a third of the terminal variance for continuous
averaging, $(n+1)(2n+1)/6n^2$ for $n$ fixings, so an at-the-money Asian is worth roughly $1/\sqrt3\approx58\%$ of the vanilla
(60\% with twelve monthly fixings).
\iqlookfor{the variance reduction and the factor.}
\end{solution}

\begin{solution}{iq:dv:asians-lookbacks-cliquets-and-forward-starts:2}
Simulate paths at the fixing dates, compute the arithmetic and geometric payoffs on the same paths, and use the geometric
Asian's closed form as a control variate with a regression coefficient; add antithetic draws; report the standard
error. The variance reduction is typically a thousandfold.
\iqlookfor{the control variate and the error report.}
\end{solution}

\begin{solution}{iq:dv:asians-lookbacks-cliquets-and-forward-starts:3}
By homogeneity: at $t_1$ it is worth $S_{t_1}C_{\mathrm{BS}}(1,k,t_2-t_1)$, so today $S_0e^{-qt_1}C_{\mathrm{BS}}(1,k,t_2-t_1)$. In a
smile model it depends on the \omterm{def:dv:local-volatility:forward}{forward smile}, the model's forecast of the smile at $t_1$ for the period to $t_2$.
\iqlookfor{homogeneity, and the \omterm{def:dv:local-volatility:forward}{forward smile}.}
\end{solution}

\begin{solution}{iq:dv:asians-lookbacks-cliquets-and-forward-starts:4}
To the \omterm{def:dv:local-volatility:forward}{forward smile}: each period is a forward-start call spread less a put spread, close to a digital when the caps are
tight, so the forward skew sets its value. Two models can match every vanilla (terminal distributions) and still assign
different joint laws to the forward returns.
\iqlookfor{forward skew, and terminal against joint distributions.}
\end{solution}

\begin{solution}{iq:dv:asians-lookbacks-cliquets-and-forward-starts:5}
Its skew is a fixed function of spot and time calibrated today, and its \omterm{def:dv:local-volatility:forward}{forward smiles} flatten: it underprices forward
skew and the \omterm{def:dv:stochastic-volatility:sv}{volatility of volatility}, so it underprices \omterm{def:dv:asians-lookbacks-cliquets-and-forward-starts:cliquet}{cliquets} and \omterm{def:dv:asians-lookbacks-cliquets-and-forward-starts:reverse}{reverse cliquets} (by 45\% against Heston in the
chapter's example).
\iqlookfor{flattening \omterm{def:dv:local-volatility:forward}{forward smiles} and the direction of the error.}
\end{solution}

\begin{solution}{iq:dv:asians-lookbacks-cliquets-and-forward-starts:6}
Short forward puts over many periods: exposed to forward volatility, to the forward skew and to the \omterm{def:dv:stochastic-volatility:sv}{volatility of
volatility}, and to crashes that hit several periods at once. Hedge the \omterm{def:dv:greeks-and-the-hedging-pnl:greeks}{vega} by forward bucket with variance or forward-start
products, the skew with risk reversals, and hold reserves for the model gap and for the \omterm{def:dv:stochastic-volatility:sv}{volatility of volatility}, which has
no liquid hedge.
\iqlookfor{the exposures, bucketed hedges and the unhedgeable part.}
\end{solution}
